\chapter*{TP : Crée ton propre jeu - Attrape les objets avec Pygame}

\begin{center}

\begin{minipage}{0.8\textwidth}
\begin{Objectif}
  \begin{center}
    Comprendre le fonctionnement d’un petit jeu en Python avec Pygame.\\
    Modifier et personnaliser un jeu selon tes envies.\\
    Compléter une partie du code manquante.\\
    Ajouter de nouvelles fonctionnalités.\\
\end{center}
\end{Objectif}
\end{minipage}
\end{center}



\section*{Le jeu}
Présentons le jeu\\
% Answer: [trim={left bottom right top},clip]
    \hspace{5pt}
    \hspace{5pt}
    \\
    \begin{itemize}
    \item[$\bullet$] Le joueur (carré rouge en bas de l’écran) se déplace à gauche et à droite avec les flèches du clavier.
    \item[$\bullet$] Un objet (carré vert) tombe du haut de l’écran.
    \item[$\bullet$] Si le joueur attrape l’objet, le score augmente et un nouvel objet apparaît.
    \item[$\bullet$] Si l’objet atteint le bas de l’écran, la partie est terminée.
\end{itemize}

\begin{Avous}
    Ouvrez Edupython puis ouvrez le code "tp-ramasse-objet.py" situé dans votre répertoire de groupe. Comme d'habitude, copiez le dans votre répertoire perso et enfin
    \begin{center}
    \large{\textbf{Testez le jeu !}}
    \end{center}

\end{Avous}
\newpage
\section*{Partie 1 : Comprendre le code}

\subsection*{Se repérer dans le jeu}
Pour placer les objets dans la fenêtre, on a besoin de se repérer. Nous allons utiliser un repère orthonormé avec des coordonnées (x,y).\\
A partir de l'image au tableau, répondre aux questions suivantes :
\begin{enumerate}[resume=globalList]
    \item Où se trouve le point (0,0) de la fenêtre ?
    \item Donner les coordonnées pour chacun des quatre coins de l'objet.
    \item Donner les coordonnées pour chacun des quatre coins du joueur.
    \item A votre avis, quelle est l'unité de mesure utilisée ?
\end{enumerate}
\begin{Comprendre}
    L'objet va se déplacer vers le bas à une vitesse donnée et le joueur va se déplacer selon les touches "flèche gauche" et "flèche droite" du clavier.
\end{Comprendre}

\subsection*{Les variables}
\begin{Avous}
    Lisez le programme qui vous est donné. Que représentent les variables suivantes ?
\end{Avous}
\vspace{-1cm}
\begin{table}[!ht]
    \centering
    \begin{tabular}{|l|p{8cm}|}
    \hline
        Variable & Rôle \\ \hline
        joueur\_x &  \\ \hline
        objet\_x, objet\_y &  \\ \hline
        VITESSE\_OBJET &  \\ \hline
        score &  \\ \hline
        jeu\_en\_cours &  \\ \hline
    \end{tabular}
\end{table}


\begin{enumerate}[resume=globalList]
    \item A la ligne 7 on lit que des variables sont des constantes. Expliquer pourquoi ces variables sont des constantes.\\
    \item A la ligne 14, on propose d'initialiser 
    \begin{itemize}
        \item la position sur l'axe x pour le joueur et les positions en x et y pour l'objet.
        \item un score
    \end{itemize} 
Expliquer le choix des valeurs qui leur ont été affectées.

\item  Que fait-on de la ligne 24 à la ligne 27 ?
\end{enumerate}


\begin{enumerate}[resume=globalList]
    \item Quelle variable permet de déplacer le joueur ? 
    \item  Comment l'objet tombe-t-il ? Comment le jeu sait-il qu’il a été attrapé ?
    \item Quelle condition met fin au jeu ? 
\end{enumerate}

\begin{Comprendre}
\textbf{Aide :}
\begin{itemize}
    \item Le jeu vérifie s’il y a une collision en comparant leurs positions.
    \item Le score augmente lorsqu’un objet est attrapé.
    \item Le jeu s’arrête si un objet atteint le bas de l’écran.
\end{itemize}
\end{Comprendre}

\begin{Avous}
Complétons un morceau de code. Remplis les trous dans le code suivant
\end{Avous}
\begin{lstlisting}[style=pythonstyle, language=python]
# Position de départ de l'objet
objet_x = __________  # Position horizontale aléatoire entre 0 et (LARGEUR - TAILLE_OBJET)
objet_y = __________  # L'objet commence en haut (0)

# L'objet tombe
objet_y += __________  # On ajoute la vitesse de chute

# Vérification de la collision avec le joueur
if __________:  # Condition pour savoir si l'objet est attrapé
    score += 1      # On augmente le score
    objet_y = 0     # L'objet revient en haut
\end{lstlisting}




\subsection*{Les fonctions}
\begin{Avous}
Changer le paramètre de la fonction delay et observer ce qui se passe.
\end{Avous}

\begin{enumerate}[resume=globalList]
\item Rappel - A quoi reconnaît-on une fonction ?
\item  A votre avis, que font chacune des fonctions suivantes ?
\end{enumerate}
\begin{table}[!ht]
    \centering
    \begin{tabular}{|l|l|p{8cm}|}
    \hline
        Fonction & Appel & Rôle \\ \hline
        init & ligne 5 & \\ \hline
        randint & ligne 16 & \\ \hline
        set\_mode & ligne 21 & \\ \hline
        set\_caption & ligne 22 & \\ \hline
        delay & ligne 35 & \\ \hline
        event.get & ligne 38 & \\ \hline
        get\_pressed & ligne 43 & \\ \hline
        fill & ligne 66 & \\ \hline
        draw.rect & ligne 67  & \\ \hline
        display.update & ligne 69 & \\ \hline
        quit & ligne 71 & \\ \hline
    \end{tabular}
\end{table}
\begin{Avous}
    Changer les paramètres des fonctions de set\_mode et set\_caption et observez ce qui se passe.
\end{Avous}

Bravo, vous maîtrisez maintenant le programme qui vous a été fourni... Allez-vous pouvoir relever les défis qui vous attendent ?!

\newpage

\section*{Partie 2 : Modifier et personnaliser le jeu}
\begin{center}
\begin{minipage}{0.5\textwidth}

    \begin{Objectif}
        \centering
        Relevez un maximum de défis !
    \end{Objectif}
\end{minipage}
    \end{center}


\subsection*{Défi 1 : Change la couleur du joueur}
\begin{lstlisting}[style=pythonstyle, language=python]
COULEUR_JOUEUR = (255, 0, 0)  # Rouge
\end{lstlisting}

\subsection*{Défi 2 : Augmente la vitesse de l’objet et observe ce qui change.}
\begin{lstlisting}[style=pythonstyle, language=python]
VITESSE_OBJET = 6
\end{lstlisting}

\subsection*{Défi 3 : Agrandis la taille du joueur et des objets.}
\begin{lstlisting}[style=pythonstyle, language=python]
TAILLE_JOUEUR = 120
TAILLE_OBJET = 40
\end{lstlisting}


\subsection*{Défi 4 : Modifie le texte affiché à la fin du jeu.}
\begin{lstlisting}[style=pythonstyle, language=python]
if objet_y > HAUTEUR:
    print(f"Perdu! Votre score est : {score}")
    jeu_en_cours = False
\end{lstlisting}
\vspace{1cm}
\section*{Partie 3 : Afficher le score en temps réel}

\subsection*{Défi 5: Ajoute l'affichage du score sur l'écran}

\subsubsection*{Ajoute une police d'écriture pour afficher le score}

Une police d'écriture est un style de texte utilisé pour afficher des informations à l’écran. Elle permet de choisir la taille et l’apparence des lettres. Dans Pygame, nous devons créer une police avant d’afficher du texte sur la fenêtre du jeu.

La police doit être ajoutée au début du programme, avant de l'utiliser pour afficher le score.
\begin{lstlisting}[style=pythonstyle, language=python]
police_score = pygame.font.Font(None, 36)
\end{lstlisting}

\subsubsection{Crée une fonction pour afficher le score}
\begin{lstlisting}[style=pythonstyle, language=python]
def afficher_score():
   texte_score = police_score.render(f"Score : {score}", True, (255, 255, 255))
   fenetre.blit(texte_score, (10, 10))
\end{lstlisting}

\subsubsection{Appelle cette fonction dans la boucle principale}
\begin{lstlisting}[style=pythonstyle, language=python]
afficher_score()
\end{lstlisting}


\newpage
\section*{Partie 4 : Rendre le jeu plus difficile}
\subsection*{Défi 6: Augmenter la difficulté du jeu}
\noindent Ajoute une variable de difficulté
\begin{lstlisting}[style=pythonstyle, language=python]
facteur_vitesse = 1.0
\end{lstlisting}
Modifie la vitesse de chute de l’objet
\begin{lstlisting}[style=pythonstyle, language=python]
objet_y += VITESSE_OBJET * facteur_vitesse
\end{lstlisting}

\subsection*{Défi 7 : Augmente la difficulté quand un objet est attrapé}
S'il y a une collision
\begin{lstlisting}[style=pythonstyle, language=python]
   facteur_vitesse += 0.1
\end{lstlisting}

\vspace{0.8cm}
\section*{Partie 5 : Ajouter de nouvelles fonctionnalités}

\subsection*{Défi 8 : Ajoute des objets spéciaux}
Ajoute ces deux morceaux de codes aux bons endroits
\begin{lstlisting}[style=pythonstyle, language=python]
objets ={'x': random.randint(0, LARGEUR - TAILLE_OBJET),
           'y': 0,
           'type': random.choice(['normal', 'bonus', 'malus'])}
\end{lstlisting}
\begin{lstlisting}[style=pythonstyle, language=python]
if objet['type'] == 'bonus':
    score += 5
elif objet['type'] == 'malus':
    score -= 3
else:
    score += 1
\end{lstlisting}

\subsection*{Défi 9 : Ajoute plusieurs objets simultanément}

\begin{lstlisting}[style=pythonstyle, language=python]
objets = [{'x': random.randint(0, LARGEUR - TAILLE_OBJET),
           'y': 0,
           'type': random.choice(['normal', 'bonus', 'malus'])}
         for _ in range(3)]
\end{lstlisting}

\subsection*{Défi 10 : Utilise une image plutôt qu'un carré pour représenter le joueur et l'objet}
\begin{lstlisting}[style=pythonstyle, language=python]
image_joueur = pygame.image.load("chemin/vers/joueur.png")
image_objet = pygame.image.load("chemin/vers/objet.png")
\end{lstlisting}

\vspace{1cm}
\begin{center}
    \huge{ Tu as relévé tous les défis ?!}\\
    \huge{Bravo }
\end{center}

Tu as appris à analyser, modifier et améliorer un jeu existant avec Pygame.  Continue à tester et à améliorer ton jeu en ajoutant des nouvelles idées !

% \newpage

% \huge{\subsection*{Des idées pour vous aider !}}

% \subsubsection*{Idée 1 : Faire tomber plusieurs objets}
% Par exemple : À éviter (fait perdre), Bonus (vaut 5 points), Booster (augmente la vitesse du joueur), Champignon (augmente la taille du joueur)...

% \subsubsection*{Idée 2 : Donner un type aux objets}
% \subsubsection*{Idée 3 : Faire sauter le joueur}
% \subsection*{Idée 4 : Transformer les objets en images}
% \subsection{}
